% ============================================================
%  VIPGuard Adversarial Robustness Extension — Full Report
%  Compile with: pdflatex -shell-escape vipguard_report.tex
%  (run twice to resolve cross-references)
% ============================================================
\documentclass[12pt,a4paper]{article}

% ── Packages ──────────────────────────────────────────────────────────────────
\usepackage[margin=2.5cm]{geometry}
\usepackage{amsmath,amssymb,amsthm}
\usepackage{graphicx}
\usepackage{booktabs}
\usepackage{array}
\usepackage{xcolor}
\usepackage{hyperref}
\usepackage{listings}
\usepackage{caption}
\usepackage{subcaption}
\usepackage{float}
\usepackage{enumitem}
\usepackage{parskip}
\usepackage{fancyhdr}
\usepackage{titlesec}
\usepackage{tcolorbox}
\usepackage{mdframed}
\usepackage{microtype}
\usepackage{multirow}
\usepackage{longtable}
\usepackage{tabularx}
\usepackage{boldline}

% ── Colour Scheme ─────────────────────────────────────────────────────────────
\definecolor{deepblue}{RGB}{0,70,140}
\definecolor{codegreen}{RGB}{0,128,64}
\definecolor{codegray}{RGB}{128,128,128}
\definecolor{codepurple}{RGB}{128,0,128}
\definecolor{codebackground}{RGB}{248,248,248}
\definecolor{alertred}{RGB}{180,0,0}
\definecolor{notebox}{RGB}{240,248,255}
\definecolor{noteborder}{RGB}{100,149,237}

% ── Hyperref Setup ────────────────────────────────────────────────────────────
\hypersetup{
    colorlinks=true,
    linkcolor=deepblue,
    citecolor=deepblue,
    urlcolor=deepblue,
    pdftitle={VIPGuard Adversarial Robustness Extension},
    pdfauthor={VIPGuard Research},
    pdfsubject={Adversarial Machine Learning},
}

% ── Code Listing Style ────────────────────────────────────────────────────────
\lstdefinestyle{pythonstyle}{
    language=Python,
    backgroundcolor=\color{codebackground},
    commentstyle=\color{codegreen}\itshape,
    keywordstyle=\color{deepblue}\bfseries,
    stringstyle=\color{codepurple},
    numberstyle=\tiny\color{codegray},
    basicstyle=\ttfamily\small,
    breakatwhitespace=false,
    breaklines=true,
    captionpos=b,
    keepspaces=true,
    numbers=left,
    numbersep=8pt,
    showspaces=false,
    showstringspaces=false,
    showtabs=false,
    tabsize=4,
    frame=single,
    rulecolor=\color{codegray},
    xleftmargin=15pt,
    xrightmargin=5pt,
    aboveskip=8pt,
    belowskip=8pt,
}
\lstset{style=pythonstyle}

% ── Note / Warning Boxes ──────────────────────────────────────────────────────
\newmdenv[
    backgroundcolor=notebox,
    linecolor=noteborder,
    linewidth=1.5pt,
    leftmargin=0pt,
    rightmargin=0pt,
    innerleftmargin=10pt,
    innerrightmargin=10pt,
    innertopmargin=8pt,
    innerbottommargin=8pt,
]{notebox}

\newmdenv[
    backgroundcolor=red!5,
    linecolor=alertred,
    linewidth=1.5pt,
    leftmargin=0pt,
    rightmargin=0pt,
    innerleftmargin=10pt,
    innerrightmargin=10pt,
    innertopmargin=8pt,
    innerbottommargin=8pt,
]{warnbox}

% ── Header / Footer ───────────────────────────────────────────────────────────
\pagestyle{fancy}
\fancyhf{}
\rhead{\textcolor{codegray}{\small VIPGuard Adversarial Robustness Report}}
\lhead{\textcolor{codegray}{\small \leftmark}}
\cfoot{\thepage}
\renewcommand{\headrulewidth}{0.4pt}

% ── Section Formatting ────────────────────────────────────────────────────────
\titleformat{\section}{\large\bfseries\color{deepblue}}{\thesection.}{0.6em}{}
\titleformat{\subsection}{\normalsize\bfseries\color{deepblue}}{\thesubsection.}{0.6em}{}
\titleformat{\subsubsection}{\normalsize\bfseries}{\thesubsubsection.}{0.6em}{}

% ── Maths helpers ─────────────────────────────────────────────────────────────
\newcommand{\norm}[1]{\left\lVert#1\right\rVert}
\newcommand{\sign}{\operatorname{sign}}

% ══════════════════════════════════════════════════════════════════════════════
\begin{document}

% ── Title Page ────────────────────────────────────────────────────────────────
\begin{titlepage}
\centering
\vspace*{2cm}
{\Huge\bfseries\color{deepblue} VIPGuard}\par
\vspace{0.4cm}
{\Large\bfseries Adversarial Robustness Extension}\par
\vspace{0.4cm}
{\large Complete Technical Report}\par
\vspace{2cm}
\rule{0.8\linewidth}{0.5pt}\par
\vspace{0.5cm}
{\large
\begin{tabular}{ll}
\textbf{System:}  & VIP Face Verification and Deepfake Detection \\
\textbf{Extension:} & FGSM / PGD Adversarial Attack \& Defence \\
\textbf{Base Model:} & Qwen2.5-VL (32B, frozen) + TransFace ViT \\
\textbf{Trainable:} & VIP Token Embeddings (32 $\times$ 3584 = 114,688 params) \\
\end{tabular}
}\par
\vspace{0.5cm}
\rule{0.8\linewidth}{0.5pt}\par
\vspace{2cm}
{\large April 2026}\par
\vfill
\end{titlepage}

% ── Table of Contents ─────────────────────────────────────────────────────────
\tableofcontents
\newpage

% ══════════════════════════════════════════════════════════════════════════════
\section{Introduction}
% ══════════════════════════════════════════════════════════════════════════════

\subsection{What is VIPGuard?}

VIPGuard is a deep-learning system that answers one question: \textit{``Is the person in this image a registered VIP, or is it a deepfake or an impostor?''} It was designed to protect against two threats simultaneously:

\begin{itemize}
    \item \textbf{Deepfake attacks}: A bad actor creates a fake video or image that makes another person's face look like the VIP's face. Ordinary face recognition systems can be fooled by high-quality deepfakes.
    \item \textbf{Impostor attacks}: A different real person tries to pass as the VIP.
\end{itemize}

Unlike a traditional face recognition system that only outputs a similarity score, VIPGuard uses a large Vision-Language Model (VLM) to reason about the image and explain its decision in natural language. This makes the system more interpretable and more capable of handling nuanced cases.

\textbf{What makes VIPGuard different from regular face recognition?}

\begin{enumerate}
    \item It produces a full reasoning chain, not just a number.
    \item It uses learned VIP-specific ``identity tokens'' that encode what a particular VIP looks and appears like.
    \item It combines two sources of information: a face similarity score from a dedicated face embedding model (TransFace) and deep visual reasoning from a 32-billion-parameter Vision-Language Model (Qwen2.5-VL).
\end{enumerate}

A typical output from VIPGuard looks like:
\begin{lstlisting}[language={},caption={Example VIPGuard output},frame=single]
<Conclusion>[Yes] The two images are of the same person.</Conclusion>
<Analysis>
1. Outer face region: [No] The jaw shape matches the VIP profile.
2. Eye region: [No] Eye spacing is consistent with registered profile.
3. Nose region: [No] Nose bridge shape matches.
4. Skin region: [No] Skin texture is consistent.
Overall Conclusion: This is the registered VIP.
</Analysis>
\end{lstlisting}

\subsection{The Problem This Report Addresses: Adversarial Robustness}

Even though VIPGuard was excellent at detecting deepfakes and impostors on normal images, it had a critical vulnerability: \textbf{adversarial attacks}.

An adversarial attack is a carefully engineered, tiny modification to an image that is invisible to the human eye but drastically changes what a machine learning model ``sees''. In the context of VIPGuard:

\begin{itemize}
    \item An attacker takes a \textit{real photo of the VIP}.
    \item They add tiny pixel-level noise (within $\varepsilon = 8/255 \approx 0.031$ of the maximum pixel value).
    \item This noise is crafted to specifically fool the face embedding model.
    \item The resulting image looks completely normal to a human but gets a very low similarity score from TransFace.
    \item VIPGuard then says ``No, this is not the VIP'' --- which is wrong.
\end{itemize}

This is the \textbf{adversarial evasion} problem: an attacker can make the system reject the real VIP's face.

\subsection{Goals of This Extended Work}

This project extended VIPGuard with a complete adversarial robustness pipeline:

\begin{enumerate}
    \item \textbf{Implement FGSM and PGD attacks} against the TransFace face embedding model.
    \item \textbf{Generate an adversarial dataset} covering all 22 VIP identities.
    \item \textbf{Build a robustness evaluation framework} measuring accuracy across four categories: normal images, deepfakes, FGSM attacks, and PGD attacks.
    \item \textbf{Design and run adversarial training} experiments to make VIPGuard robust against these attacks.
    \item \textbf{Document all design decisions, failures, and lessons learned} for future research.
\end{enumerate}

% ══════════════════════════════════════════════════════════════════════════════
\section{Original VIPGuard Architecture}
% ══════════════════════════════════════════════════════════════════════════════

\subsection{High-Level Overview}

VIPGuard has two main sub-systems working together:

\begin{enumerate}
    \item \textbf{TransFace} --- A Vision Transformer trained specifically for face recognition. It converts a face image into a 512-dimensional numerical vector called an \textit{embedding}. Two images of the same person should have embeddings that are very close to each other (high cosine similarity).
    \item \textbf{Qwen2.5-VL} --- A 32-billion-parameter Vision-Language Model that can process both images and text. It has been adapted to include \textit{VIP tokens}: small learnable embeddings that encode the identity of a specific VIP person.
\end{enumerate}

\subsection{Step-by-Step Inference Pipeline}

Here is the complete journey from input image to final decision:

\subsubsection{Step 1: Face Similarity Scoring (TransFace)}

\begin{lstlisting}[caption={Similarity scoring in inference pipeline}]
# Load the image
img = cv2.imread(image_path)           # BGR, any size

# Resize to TransFace's expected size
img = cv2.resize(img, (112, 112))

# Convert BGR to RGB and reshape to (C, H, W)
img = cv2.cvtColor(img, cv2.COLOR_BGR2RGB).transpose(2, 0, 1)

# Normalise pixels from [0, 255] to [-1, 1]
x = torch.tensor(img, dtype=torch.float32).unsqueeze(0) / 255.0
x = (x - 0.5) / 0.5   # Result range: [-1.0, 1.0]

# Forward pass through TransFace ViT
# NOTE: returns a 3-tuple; only the first element is the embedding
feat, weight, patch_entropy = fg_face.face_model(x)
# feat shape: (1, 512)

# Load the pre-computed VIP center embedding
center = torch.load('FaceDATA/FaceEmb_Center/id0.ckpt')
# center shape: (512,) -- a raw tensor, NOT a dictionary

# Compute cosine similarity and map to [0, 100]
emb_norm   = F.normalize(feat.squeeze(0), dim=0)   # unit vector
center_norm = F.normalize(center, dim=0)            # unit vector
cos_sim     = torch.dot(emb_norm, center_norm)      # range [-1, 1]
score       = int((0.5 + 0.5 * cos_sim) * 100)     # range [0, 100]
\end{lstlisting}

The formula $\text{score} = \lfloor (0.5 + 0.5 \cdot \cos(\mathbf{e}, \mathbf{c})) \times 100 \rfloor$ maps the cosine similarity (which ranges from $-1$ to $+1$) to an integer between 0 and 100. A score of 100 means perfect match, 50 means orthogonal (random chance), and 0 means opposite embeddings.

\textbf{In practice for real VIPs}: scores of 85--95/100 are typical.\newline
\textbf{For deepfakes}: scores of 52--78/100 are typical.\newline
\textbf{For FGSM adversarial images}: scores of 47--75/100.\newline
\textbf{For PGD adversarial images}: scores of 14--34/100.

\subsubsection{Step 2: VIPGuard VLM Inference}

The VLM prompt contains the similarity score as text and the image:

\begin{lstlisting}[caption={VIPGuard prompt format}]
prompt = (
    "<|face_pad|>Please determine whether the person in the input "
    "image is VIP user. "
    f"The face similarity between the input face and VIP user is "
    f"{score}/100. "
    "The face tokens are shown as follows, <|face_pad|> . "
    "You should first give your answer by 'yes' or 'no'. "
    "Then, you should explain your reasoning step by step based on "
    "different facial attributes."
)
\end{lstlisting}

Notice two things:
\begin{enumerate}
    \item The \texttt{<|face\_pad|>} tokens appear twice in the prompt. These are special placeholder tokens where the VIP identity embeddings will be inserted.
    \item The numerical similarity score is embedded in the text itself, giving the LLM an explicit numerical hint.
\end{enumerate}

\subsubsection{Step 3: The FaceChecker Module}

The FaceChecker is the bridge between the VIP token embeddings and the LLM. It is a stack of CrossAttention and FeedForward blocks.

\begin{lstlisting}[caption={FaceChecker.py -- CrossAttention and FaceChecker classes (condensed)}]
class CrossAttention(nn.Module):
    def __init__(self, embed_dim=256, num_heads=8, dropout=0.1):
        super().__init__()
        # Three linear projections: query, key, value
        self.q_proj  = nn.Linear(embed_dim, embed_dim)
        self.k_proj  = nn.Linear(embed_dim, embed_dim)
        self.v_proj  = nn.Linear(embed_dim, embed_dim)
        self.out_proj = nn.Linear(embed_dim, embed_dim)
        self.dropout = nn.Dropout(dropout)

    def forward(self, query, key, value, mask=None):
        B, Q_len, _ = query.size()
        _, KV_len, _ = key.size()

        # Project all three inputs
        Q = self.q_proj(query)    # (B, Q_len, embed_dim)
        K = self.k_proj(key)      # (B, KV_len, embed_dim)
        V = self.v_proj(value)    # (B, KV_len, embed_dim)

        # Split into multiple attention heads
        Q = Q.view(B, Q_len, self.num_heads, self.head_dim).transpose(1,2)
        K = K.view(B, KV_len, self.num_heads, self.head_dim).transpose(1,2)
        V = V.view(B, KV_len, self.num_heads, self.head_dim).transpose(1,2)

        # Scaled dot-product attention
        attn = torch.matmul(Q, K.transpose(-2, -1)) / (self.head_dim**0.5)
        attn = F.softmax(attn, dim=-1)
        out  = torch.matmul(self.dropout(attn), V)

        # Merge heads and project output
        out = out.transpose(1,2).contiguous().view(B, Q_len, self.embed_dim)
        return self.out_proj(out)


class FaceChecker(nn.Module):
    def __init__(self, in_dim, depth, num_heads=8, out_dim=3584):
        super().__init__()
        self.blocks = nn.ModuleList()
        for i in range(depth):
            if i != depth - 1:
                self.blocks.append(CrossAttention(in_dim, num_heads))
                self.blocks.append(FeedForward(in_dim, 2*in_dim, in_dim))
            else:
                # Last block projects to LLM hidden dimension (3584)
                self.blocks.append(CrossAttention(in_dim, num_heads))
                self.blocks.append(FeedForward(in_dim, 2*in_dim, out_dim))

        # VIP prompt: learnable identity tokens, shape (N, in_dim)
        # N is set by the Wrapper (default 32)
        self.vip_prompt = torch.randn(506, in_dim)

    def forward(self, img_feature_q, img_feature_v):
        # img_feature_q: image features from Qwen2.5-VL visual encoder (query)
        # img_feature_v: VIP token embeddings (key and value)
        feature = img_feature_v
        for block in self.blocks:
            if isinstance(block, CrossAttention):
                # Image features QUERY the VIP tokens
                feature = block(query=img_feature_q,
                                key=feature, value=feature)
            elif isinstance(block, FeedForward):
                feature = block(feature)
        return feature
\end{lstlisting}

\textbf{Intuition}: The image features ``ask questions'' about the VIP tokens. The VIP tokens ``answer'' based on what identity they encode. The result is a face embedding that tells the LLM how strongly this image matches the registered VIP identity.

\subsubsection{Step 4: VIP Token Injection into LLM}

The output of FaceChecker replaces the \texttt{<|face\_pad|>} positions in the input sequence. The 32-billion-parameter LLM then processes:
\begin{itemize}
    \item The face embeddings from FaceChecker at the \texttt{<|face\_pad|>} positions.
    \item The image patch tokens from the visual encoder.
    \item The text tokens (the prompt with the similarity score).
\end{itemize}

All 80 transformer layers of the LLM are \textbf{frozen} --- they are not updated during training. Only the VIP token embeddings (32 $\times$ 3584 = 114,688 parameters) are trainable.

\subsection{The Wrapper Class}

The \texttt{Wrapper} class manages all VIP-related operations:

\begin{lstlisting}[caption={Models/Wrapper.py -- Full listing}]
class Wrapper(nn.Module):
    def __init__(self, vl_model, processor, optim_facechecker=False,
                 vip_token_num=196):
        super().__init__()
        self.vl_model = vl_model

        # FREEZE the entire 32B VLM immediately
        self.vl_model.requires_grad_(False)

        self.processor = processor
        self.tokenizer = processor.tokenizer

        # Read the token dimension from FaceChecker's current prompt
        vip_token_dim = self.vl_model.facechecker.face_checker \
                                                  .vip_prompt.shape[1]
        self.optim_facechecker = optim_facechecker

        # Create a NEW learnable parameter for VIP tokens
        # This REPLACES the static vip_prompt in FaceChecker
        param = torch.empty([vip_token_num, vip_token_dim], device='cuda')
        param = nn.Parameter(nn.init.normal_(param, mean=0, std=1))
        self.vl_model.facechecker.face_checker.vip_prompt = param

    def get_optim_params(self):
        """Return only the trainable parameters."""
        optim_params = []
        for n, p in self.vl_model.named_parameters():
            if 'vip_prompt' in n:   # only VIP tokens by default
                optim_params.append(p)
        return optim_params

    def save_vip(self):
        """Save VIP token checkpoint as a dict."""
        return {
            'vip_token': self.vl_model.facechecker.face_checker.vip_prompt
        }

    def forward(self, **kwargs):
        return self.vl_model.forward(**kwargs)

    def generate(self, **kwargs):
        return self.vl_model.generate(**kwargs)
\end{lstlisting}

\textbf{Key point}: When the \texttt{Wrapper} initialises, it replaces the static \texttt{vip\_prompt} tensor inside \texttt{FaceChecker} with a proper \texttt{nn.Parameter}. This is critical because PyTorch only tracks parameters declared as \texttt{nn.Parameter} in the optimiser.

\subsection{VIP Token Checkpoint Format}

\begin{lstlisting}[caption={How VIP tokens are saved and loaded}]
# Saving (inside save_vip())
state = {'vip_token': tensor_of_shape_[32, 3584]}
torch.save(state, 'checkpoints/Stage3/id0_train/vip_token.pt')

# Loading
state = torch.load('vip_token.pt', map_location='cpu')
target = model.vl_model.facechecker.face_checker.vip_prompt
target.data.copy_(state['vip_token'])
\end{lstlisting}

\subsection{Training Stages}

The original VIPGuard was trained in two stages before this work:

\subsubsection{Stage 1: Visual Encoder Pretraining}
The visual encoder (Qwen2.5-VL's ViT) was pretrained on general image-text tasks. This stage used the standard ms-swift training framework and produced the Stage 2 checkpoint.

\subsubsection{Stage 2: Facial Attribute Training}
The full 32B model was fine-tuned on facial attribute understanding tasks using a large dataset of facial attribute questions and answers. This produced the base checkpoint at \texttt{checkpoints/checkpoints\_attr\_Stage2\_merge/}.

\subsubsection{Stage 3: Per-VIP Token Fine-tuning}
For each of the 22 VIP identities, a small set of VIP tokens (32 $\times$ 3584) is fine-tuned on that VIP's images. The rest of the model is frozen. This is the stage this work extends with adversarial training.

% ══════════════════════════════════════════════════════════════════════════════
\section{Dataset Structure}
% ══════════════════════════════════════════════════════════════════════════════

\subsection{VIP Identities}

The dataset contains 22 VIP identities: \texttt{id0} through \texttt{id21}. Each identity has its own:
\begin{itemize}
    \item Training images organised into pair-type subdirectories.
    \item A pre-computed face embedding centre (\texttt{FaceEmb\_Center/id\{k\}.ckpt}).
    \item A pre-trained VIP token file (\texttt{Pretrained\_VIPToken/id\{k\}.pt}).
    \item A training JSON file (\texttt{Stage3\_training\_json/id\{k\}.json}).
\end{itemize}

\subsection{Image Pair Categories}

\begin{table}[H]
\centering
\caption{Image pair categories and their meanings}
\begin{tabular}{@{}llll@{}}
\toprule
\textbf{Directory} & \textbf{Full Name} & \textbf{Label} & \textbf{Description} \\
\midrule
\texttt{r\_r\_i}   & Real vs Real, same Identity & 0 (Yes) & Real VIP image \\
\texttt{r\_r\_d\_i} & Real vs Real, Different Identity & 1 (No) & Different person \\
\texttt{r\_fs\_i}  & Real vs Face-Swapped & 1 (No) & Deepfake (face swap) \\
\texttt{r\_efs\_i} & Real vs Entire-Face-Synthesised & 1 (No) & Deepfake (full synthesis) \\
\texttt{adv\_real} & Adversarial Real & 0 (Yes) & Perturbed real VIP (new) \\
\bottomrule
\end{tabular}
\end{table}

\textbf{Note on label convention}: Label 0 means \textit{the model should say ``Yes'' (this is the VIP)}. Label 1 means \textit{the model should say ``No'' (this is not the VIP)}. This is slightly counterintuitive but matches the PyTorch convention used in the codebase.

\subsection{Folder Structure}

\begin{lstlisting}[language={},caption={Repository directory structure}]
VIPGuard/
|-- FaceDATA/
|   |-- Training_Img/
|   |   |-- id0/
|   |   |   |-- r_r_i/          # Real same-identity pairs
|   |   |   |   |-- id0_38/
|   |   |   |       |-- 2.png   # Probe image (used for attacks)
|   |   |   |       |-- real.png
|   |   |   |-- r_r_d_i/        # Different-identity pairs
|   |   |   |-- r_fs_i/         # Face-swapped deepfakes
|   |   |   |   |-- id0_12_face/
|   |   |   |       |-- fake.png
|   |   |   |-- r_efs_i/        # Entire-face-synthesis deepfakes
|   |   |-- id1/ ... id21/
|   |-- FaceEmb_Center/
|   |   |-- id0.ckpt ... id21.ckpt  # Raw (512,) tensors
|   |-- Pretrained_VIPToken/
|   |   |-- id0.pt ... id21.pt      # {'vip_token': tensor([N, 3584])}
|   |-- Stage3_training_json/
|       |-- id0.json ... id21.json
|
|-- data/adversarial/            # Generated by this work
|   |-- fgsm/adv_real/id{k}/{pair}/
|   |   |-- adv.png              # FGSM adversarial image
|   |   |-- sim.txt              # Adversarial similarity score
|   |   |-- orig_sim.txt         # Original pre-attack similarity score
|   |-- pgd/adv_real/id{k}/{pair}/  # Same structure for PGD
|   |-- json/
|       |-- id{k}_adv_mixed.json    # Mixed training JSON
|
|-- checkpoints/
|   |-- checkpoints_attr_Stage2_merge/  # Base 32B model (frozen)
|   |-- Stage3_scratch/                 # Adversarially trained tokens (new)
|       |-- id0/vip_token.pt
|
|-- results/
|   |-- before_adv_train/       # Baseline results
|   |-- after_scratch_train/    # Post-training results
|   |-- all_runs_summary.json   # All experimental results
|
|-- attacks/                    # New: attack modules
|   |-- __init__.py
|   |-- fgsm.py
|   |-- pgd.py
|   |-- utils.py
|
|-- Models/
|   |-- Wrapper.py
|   |-- Face_Model/
|   |   |-- FaceModel.py
|   |   |-- TransFace/backbones/vit.py
|   |-- VIPGuard/
|       |-- FaceChecker.py
|       |-- modeling_qwen2_5_vl.py
|       |-- processing_qwen2_5_vl.py
|
|-- Stage3_train.py
|-- VIP_Dataset.py
|-- generate_adversarial_dataset.py
|-- evaluate_robustness.py
|-- adversarial_training.py
|-- scratch_adv_training.py
|-- monitor_pipeline.py
|-- CLAUDE.md
\end{lstlisting}

\subsection{VIP\_Dataset.py --- Dataset Loader}

This class loads the training JSON and prepares batches for the training loop:

\begin{lstlisting}[caption={VIP\_Dataset.py -- Key methods}]
class VIP_Dataset(Dataset):
    def __init__(self, json_path, processor):
        self.json_file = self.read_json(json_path)  # Load all entries
        self.processor = processor

    def load_data(self, json_data):
        prompt  = json_data['messages'][0]['content']
        img_dir = json_data['images'][0]    # Path to image file
        response = json_data['messages'][1]['content']  # Expected answer

        # Build the chat structure expected by Qwen2.5-VL processor
        chat = [
            {'role': 'user', 'content': [
                {'type': 'image', 'image': img_dir},
                {'type': 'text',  'text': prompt}
            ]},
            {'role': 'assistant', 'content': [
                {'type': 'text', 'text': response}
            ]}
        ]

        # CRITICAL: Label assignment by path substring
        # 'r_r_i'   --> label 0 (authentic VIP)
        # 'adv_real' --> label 0 (adversarially perturbed VIP)
        # everything else --> label 1 (impostor / deepfake)
        if 'r_r_i' in img_dir.lower() or 'adv_real' in img_dir.lower():
            label = 0
        else:
            label = 1

        return question, answer, answer_empty, img_dir, label
\end{lstlisting}

\begin{warnbox}
\textbf{Important patch}: The original code only checked for \texttt{r\_r\_i} to assign label 0. We added \texttt{or `adv\_real' in img\_dir.lower()} so that adversarially perturbed real VIP images also get the correct label 0 (authentic). Without this patch, the adversarial training data would have the wrong labels.
\end{warnbox}

% ══════════════════════════════════════════════════════════════════════════════
\section{Adversarial Attack Implementation}
% ══════════════════════════════════════════════════════════════════════════════

\subsection{Threat Model}

\begin{itemize}
    \item \textbf{Attacker's goal}: Make a real VIP image be rejected by the system (evasion attack / dodge attack).
    \item \textbf{Knowledge}: White-box access to TransFace --- the attacker knows the model architecture and weights.
    \item \textbf{Constraint}: The perturbation must be imperceptible. We use $\varepsilon = 8/255$ in pixel space, which is a standard benchmark perturbation that humans cannot detect.
    \item \textbf{Attack target}: TransFace embeddings, not the full 32B VLM. Attacking the VLM directly would require computing gradients through 32 billion parameters, which is impractical.
\end{itemize}

\subsection{Image Normalisation Convention}

TransFace expects images normalised to $[-1, 1]$:
\begin{equation}
x_{\text{norm}} = \frac{x_{\text{pixel}} / 255 - 0.5}{0.5}
\end{equation}

This means an $\varepsilon = 8/255$ perturbation in pixel space corresponds to:
\begin{equation}
\varepsilon_{\text{norm}} = \frac{\varepsilon_{\text{pixel}}}{0.5} = \frac{8/255}{0.5} = \frac{16}{255} \approx 0.0627
\end{equation}

in normalised space. This conversion is critical --- failing to convert would either under- or over-perturb the image.

\subsection{FGSM Attack}

\subsubsection{Mathematical Background}

FGSM (Fast Gradient Sign Method) was introduced by Goodfellow et al.\ (2015). It asks: \textit{in which direction should we push each pixel to maximally hurt the model?}

The answer is given by the gradient of the loss with respect to the input image. Since we want to \textit{decrease} the similarity (push the embedding away from the VIP centre), we use gradient \textit{descent}:

\begin{equation}
\mathbf{x}_{\text{adv}} = \mathbf{x} - \varepsilon \cdot \sign\!\left(\nabla_{\mathbf{x}}\, \cos\!\left(f(\mathbf{x}),\, \mathbf{c}\right)\right)
\end{equation}

where:
\begin{itemize}
    \item $\mathbf{x}$ is the original normalised image tensor.
    \item $f(\mathbf{x})$ is the TransFace embedding of $\mathbf{x}$.
    \item $\mathbf{c}$ is the VIP centre embedding.
    \item $\cos(\cdot,\cdot)$ is cosine similarity.
    \item $\varepsilon$ is the perturbation budget.
    \item $\sign(\cdot)$ takes the element-wise sign ($+1$ or $-1$).
\end{itemize}

By subtracting the sign of the gradient, we move the image in the direction that \textbf{minimises} the cosine similarity.

\subsubsection{Complete Code Walkthrough: attacks/fgsm.py}

\begin{lstlisting}[caption={attacks/fgsm.py -- Complete listing with line-by-line explanation}]
import cv2
import numpy as np
import torch
import torch.nn.functional as F


def preprocess_for_attack(img_bgr: np.ndarray, device: str) -> torch.Tensor:
    """BGR uint8 (H,W,3) --> float32 tensor (1,3,112,112) in [-1,1]."""

    # Step 1: Resize to TransFace's expected input size
    img = cv2.resize(img_bgr, (112, 112))

    # Step 2: Convert BGR (OpenCV default) to RGB (PyTorch convention)
    img = cv2.cvtColor(img, cv2.COLOR_BGR2RGB)

    # Step 3: Rearrange from (H, W, C) to (C, H, W) for PyTorch
    img = img.transpose(2, 0, 1)  # shape: (3, 112, 112)

    # Step 4: Create tensor and scale from [0,255] to [0,1]
    x = torch.tensor(img, dtype=torch.float32).unsqueeze(0) / 255.0
    # unsqueeze(0) adds batch dimension: shape (1, 3, 112, 112)

    # Step 5: Normalise from [0,1] to [-1,1] (TransFace convention)
    x = (x - 0.5) / 0.5   # equivalent to (x - mean) / std with mean=std=0.5

    # Move to GPU (or CPU)
    return x.to(device)
    # IMPORTANT: requires_grad is NOT set here; caller sets it if needed


def postprocess_from_attack(adv_tensor: torch.Tensor) -> np.ndarray:
    """float32 tensor (1,3,112,112) in [-1,1] --> BGR uint8 (112,112,3)."""

    # Reverse the normalisation: [-1,1] --> [0,1]
    adv = adv_tensor.squeeze(0).permute(1, 2, 0).detach().cpu().numpy()
    # squeeze(0): remove batch dim. permute(1,2,0): (C,H,W) -> (H,W,C)

    # Scale from [0,1] to [0,255] and clip to valid range
    adv = ((adv * 0.5 + 0.5) * 255.0).clip(0, 255).astype(np.uint8)

    # Convert RGB back to BGR for OpenCV
    return cv2.cvtColor(adv, cv2.COLOR_RGB2BGR)


def fgsm_attack(face_backbone, img_bgr, center,
                epsilon=8/255, device="cuda") -> np.ndarray:
    """
    Single-step FGSM dodge attack on TransFace.
    face_backbone: fg_face.face_model (raw ViT) -- NOT FG_Face wrapper
    """

    # Put model in eval mode:
    # - BatchNorm uses running statistics (not batch statistics)
    # - Random patch masking is disabled
    # Both are essential for single-image attack (batch size = 1)
    face_backbone.eval()

    # Normalise the VIP centre vector to unit length and move to GPU
    # detach() prevents gradients from flowing back through center
    center_norm = F.normalize(center.float().to(device), dim=0).detach()

    # Convert epsilon from pixel space to normalised space (see Section 4.2)
    eps_norm = epsilon / 0.5   # = 8/255 / 0.5 = 0.0627

    # Preprocess image and ENABLE gradient tracking on the input
    # requires_grad_(True) tells PyTorch: "I want to compute d(loss)/d(x)"
    x = preprocess_for_attack(img_bgr, device).requires_grad_(True)

    # Forward pass through TransFace ViT
    # CRITICAL: disable autocast to prevent FP16 gradient underflow.
    # The ViT internally uses autocast; disabling it here ensures gradients
    # are computed in full float32, preventing them from being zeroed out.
    with torch.amp.autocast('cuda', enabled=False):
        feat, _, _ = face_backbone(x.float())
        # feat shape: (1, 512)  -- the 512-dimensional face embedding

    # Normalise the embedding to unit length (for cosine similarity)
    emb_norm = F.normalize(feat.squeeze(0), dim=0)
    # emb_norm shape: (512,)

    # Compute cosine similarity between this image's embedding and VIP centre
    sim = torch.dot(emb_norm, center_norm)
    # sim is a scalar in [-1, 1]

    # Backward pass: compute gradient of sim w.r.t. x
    # Since sim = cos(f(x), c), this gives us ds/dx
    # We want to MINIMISE sim (dodge: push embedding away from centre)
    sim.backward()
    # After this, x.grad contains ds/dx -- shape (1, 3, 112, 112)

    with torch.no_grad():
        # The sign of each gradient component tells us which direction
        # increases similarity. We move in the OPPOSITE direction (subtract).
        grad_sign = x.grad.sign()    # shape (1, 3, 112, 112), values in {-1, +1}

        # Apply perturbation and clamp to valid normalised range
        x_adv = (x - eps_norm * grad_sign).clamp(-1.0, 1.0)
        # x_adv is the adversarial image, still in normalised float32 format

    # Convert back to BGR uint8 for saving
    return postprocess_from_attack(x_adv)
\end{lstlisting}

\subsubsection{Why We Use \texttt{face\_backbone} Directly}

A critical pitfall: \texttt{FG\_Face.face\_emb\_forward()} is decorated with \texttt{@torch.no\_grad()}, which completely disables gradient computation for everything inside it. If we called it for the attack, all gradients would be zero, and the attack would do nothing.

\begin{lstlisting}[caption={The @no\_grad decorator in FaceModel.py}]
class FG_Face:
    @torch.no_grad()          # <-- This blocks ALL gradient flow
    def face_emb_forward(self, x):
        feat, _, _ = self.face_model(x)   # face_model is the raw ViT
        return feat

# WRONG (gradients = 0):
emb = fg_face.face_emb_forward(x)   # no_grad decorator kills gradients

# CORRECT (gradients flow):
feat, _, _ = fg_face.face_model(x)  # direct call to raw ViT
\end{lstlisting}

\subsection{PGD Attack}

\subsubsection{Mathematical Background}

PGD (Projected Gradient Descent), introduced by Madry et al.\ (2018), is an iterative version of FGSM. Instead of taking one large step, it takes many small steps, projecting back into the allowed perturbation ball after each step.

\begin{equation}
\mathbf{x}^{(0)}_{\text{adv}} = \mathbf{x} + \boldsymbol{\delta}_0, \quad \boldsymbol{\delta}_0 \sim \mathcal{U}(-\varepsilon, +\varepsilon)
\end{equation}
\begin{equation}
\mathbf{x}^{(t+1)}_{\text{adv}} = \Pi_{\mathcal{B}(\mathbf{x},\varepsilon)}\!\left(\mathbf{x}^{(t)}_{\text{adv}} - \alpha \cdot \sign\!\left(\nabla_{\mathbf{x}^{(t)}} \cos\!\left(f(\mathbf{x}^{(t)}), \mathbf{c}\right)\right)\right)
\end{equation}

where:
\begin{itemize}
    \item $\alpha$ is the step size per iteration (we use $\alpha = 2/255$).
    \item $t$ ranges from 0 to $T-1$ (we use $T = 10$ steps).
    \item $\Pi_{\mathcal{B}(\mathbf{x},\varepsilon)}$ projects the result back into the $\varepsilon$-ball around the original image: $\text{clip}(\mathbf{x}_{\text{adv}} - \mathbf{x}, -\varepsilon, +\varepsilon) + \mathbf{x}$.
    \item Random start ($\boldsymbol{\delta}_0$) avoids saddle points and produces stronger attacks.
\end{itemize}

\subsubsection{Complete Code Walkthrough: attacks/pgd.py}

\begin{lstlisting}[caption={attacks/pgd.py -- Complete listing with line-by-line explanation}]
import torch
import torch.nn.functional as F
from attacks.fgsm import preprocess_for_attack, postprocess_from_attack


def pgd_attack(face_backbone, img_bgr, center,
               epsilon=8/255, alpha=2/255, num_steps=10,
               device="cuda", random_start=True) -> np.ndarray:

    face_backbone.eval()   # Disable BatchNorm training mode (see FGSM)

    # Normalise the VIP centre
    center_norm = F.normalize(center.float().to(device), dim=0).detach()

    # Convert epsilon and alpha from pixel space to normalised [-1,1] space
    eps_n   = epsilon / 0.5   # = 0.0627 (total budget)
    alpha_n = alpha   / 0.5   # = 0.0157 (per-step budget)

    # Preprocess original image -- NO gradient needed on this
    x_orig = preprocess_for_attack(img_bgr, device).detach()

    # Random start: add uniform noise within [-eps, +eps] to initialise
    # This explores different starting points to find stronger attacks
    if random_start:
        noise = torch.empty_like(x_orig).uniform_(-eps_n, eps_n)
        x_adv = (x_orig + noise).clamp(-1.0, 1.0).detach()
    else:
        x_adv = x_orig.clone().detach()

    # --- PGD Iterative Loop ---
    for _ in range(num_steps):

        # CRITICAL: Must call .detach() BEFORE .requires_grad_(True)
        # Without detach(), the computation graph from the PREVIOUS iteration
        # is kept alive, and after 10 steps, PyTorch runs out of GPU memory.
        x_adv = x_adv.detach().requires_grad_(True)

        # Forward pass in float32 (avoid FP16 gradient underflow)
        with torch.amp.autocast('cuda', enabled=False):
            feat, _, _ = face_backbone(x_adv.float())

        emb_norm = F.normalize(feat.squeeze(0), dim=0)

        # Compute cosine similarity and backpropagate
        sim = torch.dot(emb_norm, center_norm)
        sim.backward()   # Gradient: how does each pixel affect similarity?

        with torch.no_grad():
            # Take a step in the NEGATIVE gradient direction (minimise sim)
            x_adv = x_adv - alpha_n * x_adv.grad.sign()

            # Project back into epsilon-ball around ORIGINAL image
            # This ensures the total perturbation never exceeds epsilon
            delta = (x_adv - x_orig).clamp(-eps_n, eps_n)
            x_adv = (x_orig + delta).clamp(-1.0, 1.0)
            # The final clamp keeps pixel values in the valid [-1, 1] range

    return postprocess_from_attack(x_adv)
\end{lstlisting}

\subsubsection{Key Differences: FGSM vs PGD}

\begin{table}[H]
\centering
\caption{FGSM vs PGD comparison}
\begin{tabular}{@{}lll@{}}
\toprule
\textbf{Property} & \textbf{FGSM} & \textbf{PGD} \\
\midrule
Steps & 1 & 10 \\
Per-step size & $\varepsilon$ & $\alpha = \varepsilon/4$ \\
Random start & No & Yes \\
Computation time & Fast (\textless 1s) & Slow ($\sim$10s) \\
Attack strength & Weak & Strong \\
Sim drop (typical) & 85 $\to$ 47--75 & 85 $\to$ 14--34 \\
\bottomrule
\end{tabular}
\end{table}

\subsection{Shared Utilities: attacks/utils.py}

\begin{lstlisting}[caption={attacks/utils.py -- Shared helper functions}]
def load_face_fg(device='cpu'):
    """Load full FG_Face wrapper for similarity computation."""
    from Models.Face_Model.FaceModel import FG_Face
    fg_face = FG_Face(attributes=[], token_dim=3584, model_name='transface')
    fg_face.face_model.eval()
    fg_face.face_model.to(device)
    return fg_face


def load_center(id_name, center_dir='./FaceDATA/FaceEmb_Center'):
    """Load VIP embedding centre -- a raw (512,) tensor (not a dict)."""
    path = os.path.join(center_dir, f'{id_name}.ckpt')
    center = torch.load(path, map_location='cpu')
    # Safety check: some checkpoints may be saved as dicts
    if isinstance(center, dict):
        center = center.get('face_emb', next(iter(center.values())))
    return center.float()


def cosine_sim_score(emb, center):
    """Compute the [0,100] similarity score used throughout VIPGuard."""
    emb_n = F.normalize(emb.float().cpu(), dim=0)
    ctr_n = F.normalize(center.float().cpu(), dim=0)
    cos = torch.dot(emb_n, ctr_n).item()          # range [-1, 1]
    return int((0.5 + 0.5 * cos) * 100)            # range [0, 100]


def load_vip_token_into(model, vip_token_path):
    """Swap VIP tokens into a loaded model without reloading the 32B VLM."""
    state = torch.load(vip_token_path, map_location='cpu')
    target = model.vl_model.facechecker.face_checker.vip_prompt
    if isinstance(state, dict) and 'vip_token' in state:
        token = state['vip_token']
    else:
        token = state
    # Handle token count mismatch gracefully
    if token.shape[0] != target.shape[0]:
        token = token[:target.shape[0]]
    # copy_ preserves dtype and device of target
    target.data.copy_(token.to(dtype=target.dtype, device=target.device))
\end{lstlisting}

% ══════════════════════════════════════════════════════════════════════════════
\section{Adversarial Dataset Generation}
% ══════════════════════════════════════════════════════════════════════════════

\subsection{Overview}

The script \texttt{generate\_adversarial\_dataset.py} produces:
\begin{enumerate}
    \item Adversarial images: FGSM and PGD versions of every real VIP image.
    \item Similarity files: \texttt{sim.txt} (post-attack, for evaluation) and \texttt{orig\_sim.txt} (pre-attack, for training).
    \item Mixed training JSONs: combining original clean data with adversarial examples.
\end{enumerate}

\subsection{Source Image Selection}

Only \texttt{r\_r\_i} (real same-identity) images are attacked. These are the images an adversary would target: real photos of the VIP that the system should accept but the attack tries to reject.

\begin{lstlisting}[caption={Source image enumeration}]
def get_source_images(id_name, data_root='./FaceDATA/Training_Img'):
    results = []
    # Only attack r_r_i (real same-identity) images
    base = os.path.join(data_root, id_name, 'r_r_i')
    for pair_dir in sorted(os.listdir(base)):
        # Try standard probe image filenames
        for candidate in ['2.png', 'fake.png', 'probe.png']:
            img_path = os.path.join(base, pair_dir, candidate)
            if os.path.isfile(img_path):
                results.append((img_path, pair_dir))
                break
    return results
    # Typical counts: id0=400, id1=400, id2=231 images
\end{lstlisting}

\subsection{Core Generation Loop}

\begin{lstlisting}[caption={generate\_adversarial\_dataset.py -- Core generation function}]
def generate_adv_for_id(id_name, fg_face, center, attack_type,
                         output_root, epsilon, alpha, num_steps,
                         device, skip_existing):

    backbone = fg_face.face_model    # Raw ViT for gradient attacks
    source_images = get_source_images(id_name)

    out_dir = os.path.join(output_root, attack_type, 'adv_real', id_name)
    os.makedirs(out_dir, exist_ok=True)
    json_entries = []

    for img_path, pair_name in source_images:
        pair_out      = os.path.join(out_dir, pair_name)
        adv_img_path  = os.path.join(pair_out, 'adv.png')
        sim_txt_path  = os.path.join(pair_out, 'sim.txt')
        orig_sim_path = os.path.join(pair_out, 'orig_sim.txt')
        os.makedirs(pair_out, exist_ok=True)

        # If already generated, reload from disk
        if skip_existing and os.path.isfile(adv_img_path):
            adv_sim  = int(open(sim_txt_path).read())  if os.path.isfile(sim_txt_path)  else -1
            orig_sim = int(open(orig_sim_path).read()) if os.path.isfile(orig_sim_path) else adv_sim
            json_entries.append({'img_path': ..., 'adv_sim': adv_sim,
                                 'orig_sim': orig_sim})
            continue

        img_bgr = cv2.imread(img_path)

        # Record ORIGINAL similarity BEFORE attack (for reference)
        orig_sim = sim_from_array(img_bgr, fg_face, center)

        # Apply attack (FGSM or PGD)
        if attack_type == 'fgsm':
            adv_bgr = fgsm_attack(backbone, img_bgr, center,
                                   epsilon=epsilon, device=device)
        else:
            adv_bgr = pgd_attack(backbone, img_bgr, center,
                                  epsilon=epsilon, alpha=alpha,
                                  num_steps=num_steps, device=device)

        # Record ADVERSARIAL similarity AFTER attack (for evaluation)
        adv_sim = sim_from_array(adv_bgr, fg_face, center)

        # Save files
        cv2.imwrite(adv_img_path, adv_bgr)
        open(sim_txt_path,  'w').write(str(adv_sim))   # used at eval time
        open(orig_sim_path, 'w').write(str(orig_sim))  # used at train time

        json_entries.append({'img_path': rel_path,
                             'adv_sim': adv_sim,
                             'orig_sim': orig_sim})

    return json_entries
\end{lstlisting}

\subsection{Training JSON Format}

Each entry in the mixed training JSON follows this structure:

\begin{lstlisting}[language={},caption={Training JSON entry format}]
{
  "messages": [
    {
      "role": "user",
      "content": "<|face_pad|><image>Please determine whether the person
                  in the input image is VIP user.
                  The face similarity between the input face and VIP user
                  is 63/100. The face tokens are shown as follows,
                  <|face_pad|>. You should first give your answer by
                  'yes' or 'no'. Then, you should explain your reasoning
                  step by step based on different facial attributes."
    },
    {
      "role": "assistant",
      "content": "<Conclusion>[Yes] The two images are of the same person.
                  </Conclusion>\n<Analysis>..."
    }
  ],
  "images": ["./data/adversarial/fgsm/adv_real/id0/id0_38/adv.png"]
}
\end{lstlisting}

\subsection{Mixed JSON Composition}

\begin{lstlisting}[caption={Building the mixed training dataset}]
def build_mixed_json(id_name, fgsm_entries, pgd_entries,
                      original_json_path, output_path, adv_fraction=0.15):

    original = json.load(open(original_json_path))   # Clean training data
    n_original  = len(original)       # e.g. 897 for id0

    # Total adversarial samples = 15% of original size
    n_adv_total = max(1, int(n_original * adv_fraction))   # = 134

    # Split budget equally between FGSM and PGD
    n_each  = n_adv_total // 2        # = 67 each
    all_adv = fgsm_entries[:n_each] + pgd_entries[:n_each]   # 134 total

    combined = original + all_adv     # 897 + 134 = 1031 entries
    random.Random(42).shuffle(combined)   # Reproducible shuffle

    json.dump(combined, open(output_path, 'w'), indent=2, ensure_ascii=False)
    # Result: 87% clean + 13% adversarial
\end{lstlisting}

\textbf{Data composition for id0}:
\begin{itemize}
    \item 897 original entries (natural mix of real VIP, different identities, deepfakes)
    \item 67 FGSM adversarial real (label: Yes)
    \item 67 PGD adversarial real (label: Yes)
    \item \textbf{Total}: 1,031 entries
\end{itemize}

% ══════════════════════════════════════════════════════════════════════════════
\section{Robustness Evaluation Framework}
% ══════════════════════════════════════════════════════════════════════════════

\subsection{Four Evaluation Categories}

\begin{table}[H]
\centering
\caption{Evaluation categories and their properties}
\begin{tabular}{@{}lllll@{}}
\toprule
\textbf{Category} & \textbf{Image Source} & \textbf{True Label} & \textbf{Sim Range} & \textbf{Count} \\
\midrule
normal\_real  & \texttt{r\_r\_i/*/2.png}         & Yes (1) & 79--95  & $\sim$400 \\
deepfake      & \texttt{r\_fs\_i/*, r\_efs\_i/*}  & No  (0) & 52--86  & $\sim$65  \\
fgsm\_attack  & \texttt{fgsm/adv\_real/*}         & Yes (1) & 47--75  & $\sim$400 \\
pgd\_attack   & \texttt{pgd/adv\_real/*}          & Yes (1) & 14--34  & $\sim$400 \\
\bottomrule
\end{tabular}
\end{table}

\begin{warnbox}
\textbf{Critical observation}: The similarity score ranges for FGSM (47--75) and deepfake (52--86) overlap almost completely. Both categories appear to the model as ``medium similarity''. This is the fundamental challenge of this project.
\end{warnbox}

\subsection{Evaluation Metrics}

\subsubsection{Confusion Matrix}

For each image, the model predicts Yes (1) or No (0):

\begin{table}[H]
\centering
\caption{Confusion matrix for VIPGuard evaluation}
\begin{tabular}{@{}lll@{}}
\toprule
 & \textbf{Model says Yes} & \textbf{Model says No} \\
\midrule
\textbf{True label: Yes (real VIP)} & TP (True Positive) & FN (False Negative) \\
\textbf{True label: No (impostor)}  & FP (False Positive) & TN (True Negative) \\
\bottomrule
\end{tabular}
\end{table}

\subsubsection{Metric Formulas}

\begin{align}
\text{Accuracy} &= \frac{TP + TN}{N} \\[6pt]
\text{FAR (False Acceptance Rate)} &= \frac{FP}{FP + TN} \quad \text{(impostors/fakes accepted)} \\[6pt]
\text{FRR (False Rejection Rate)} &= \frac{FN}{FN + TP} \quad \text{(real VIPs rejected)} \\[6pt]
\text{Performance Drop} &= \frac{\text{acc}_{\text{normal}} - \text{acc}_{\text{adv}}}{\text{acc}_{\text{normal}}} \times 100\%
\end{align}

\begin{itemize}
    \item \textbf{FAR} is relevant for deepfake categories. A high FAR means deepfakes are being accepted as real VIPs --- dangerous!
    \item \textbf{FRR} is relevant for normal\_real and adversarial categories. A high FRR means real VIPs are being wrongly rejected.
    \item \textbf{Performance Drop} measures how much the adversarial attack degraded performance relative to baseline.
\end{itemize}

\subsection{Evaluation Pipeline: evaluate\_robustness.py}

\subsubsection{Model Loading --- Load Once, Swap Per Identity}

\begin{lstlisting}[caption={evaluate\_robustness.py -- Efficient model loading strategy}]
def load_vipguard_base(pretrain_path, token_num,
                        device_map='auto', load_in_4bit=True):
    """
    Load the 32B VLM ONCE. Then swap the VIP token for each identity.

    Loading the full model takes ~5 minutes and ~32GB VRAM.
    Without this strategy, evaluating 22 VIPs would require 22 model loads.
    """
    if load_in_4bit:
        # NF4 quantization: compresses 32B model from 64GB to ~8GB
        bnb_config = BitsAndBytesConfig(
            load_in_4bit=True,
            bnb_4bit_quant_type='nf4',          # 4-bit normal float format
            bnb_4bit_compute_dtype=torch.float16, # computations in FP16
            bnb_4bit_use_double_quant=True,      # additional compression
        )
        vl_model = Qwen2_5_VLForConditionalGeneration.from_pretrained(
            pretrain_path,
            device_map='auto',         # spread across available GPUs
            quantization_config=bnb_config,
        )

    model = Wrapper(vl_model, processor, token_num=token_num)

    # Cast only the non-quantized FaceChecker to float16
    # (quantized model parts must stay in NF4 format)
    model.vl_model.facechecker.half()

    return model, processor
\end{lstlisting}

\subsubsection{Building Evaluation Entry Lists}

\begin{lstlisting}[caption={Building the four evaluation sets for one VIP}]
def build_evaluation_entries(id_name, data_root, adv_root, fg_face, center):
    entries = {'normal_real': [], 'deepfake': [],
               'fgsm_attack': [], 'pgd_attack': []}

    # 1. Normal real images: compute similarity live
    for pair in os.listdir(os.path.join(data_root, id_name, 'r_r_i')):
        img_path = os.path.join(data_root, id_name, 'r_r_i', pair, '2.png')
        sim = sim_from_path(img_path, fg_face, center)
        entries['normal_real'].append((img_path, sim, 1))  # true label = 1 (Yes)

    # 2. Deepfakes: r_fs_i and r_efs_i
    for fake_type in ['r_fs_i', 'r_efs_i']:
        for pair in os.listdir(os.path.join(data_root, id_name, fake_type)):
            img_path = os.path.join(data_root, id_name, fake_type, pair, 'fake.png')
            sim = sim_from_path(img_path, fg_face, center)
            entries['deepfake'].append((img_path, sim, 0))  # true label = 0 (No)

    # 3. FGSM adversarial: read sim from pre-computed sim.txt
    for pair in os.listdir(os.path.join(adv_root, 'fgsm', 'adv_real', id_name)):
        img_path = os.path.join(adv_root, 'fgsm', 'adv_real', id_name, pair, 'adv.png')
        sim_txt  = os.path.join(adv_root, 'fgsm', 'adv_real', id_name, pair, 'sim.txt')
        sim = int(open(sim_txt).read()) if os.path.isfile(sim_txt) else ...
        entries['fgsm_attack'].append((img_path, sim, 1))  # true label = 1 (Yes)

    # 4. PGD adversarial: same pattern
    ...
    return entries
\end{lstlisting}

\subsubsection{Running Inference and Parsing Output}

\begin{lstlisting}[caption={Inference and prediction parsing}]
def run_inference(img_path, sim_score, model, processor, token_num):
    prompt = (
        f"<|face_pad|>Please determine whether the person in the input "
        f"image is VIP user. The face similarity between the input face "
        f"and VIP user is {sim_score}/100. The face tokens are shown as "
        f"follows, <|face_pad|> . You should first give your answer by "
        f"'yes' or 'no'. Then, you should explain your reasoning step by "
        f"step based on different facial attributes."
    )
    # Note: sim_score for adversarial images is the POST-ATTACK score
    # from sim.txt, not the original pre-attack score

    message = [{'role': 'user', 'content': [
        {'type': 'image', 'image': img_path},
        {'type': 'text',  'text': prompt}
    ]}]
    # ... process inputs, run model.generate(...) ...
    return raw_output_text


def parse_prediction(output_text):
    """
    Parse model output to Yes (1) or No (0).
    Primary: look for <Conclusion>[Yes] or <Conclusion>[No]
    Fallback: scan lines for 'yes' or 'no'
    """
    match = re.search(r'<Conclusion>\[(Yes|No)\]', output_text, re.IGNORECASE)
    if match:
        return 1 if match.group(1).lower() == 'yes' else 0
    # ...
\end{lstlisting}

\subsubsection{Metrics Computation}

\begin{lstlisting}[caption={Metrics computation in evaluate\_category()}]
TP = sum(1 for p, l in zip(predictions, labels) if p == 1 and l == 1)
TN = sum(1 for p, l in zip(predictions, labels) if p == 0 and l == 0)
FP = sum(1 for p, l in zip(predictions, labels) if p == 1 and l == 0)
FN = sum(1 for p, l in zip(predictions, labels) if p == 0 and l == 1)

accuracy = (TP + TN) / N
FAR      = FP / (FP + TN)   # impostors accepted (dangerous)
FRR      = FN / (FN + TP)   # real VIPs rejected (annoying but safe)
\end{lstlisting}

% ══════════════════════════════════════════════════════════════════════════════
\section{Experimental Results and Analysis}
% ══════════════════════════════════════════════════════════════════════════════

\subsection{Baseline Results (No Adversarial Training)}

Before any adversarial training, the pretrained VIP tokens worked perfectly on clean images but were completely vulnerable to attacks:

\begin{table}[H]
\centering
\caption{Baseline results --- Pretrained VIP tokens, no adversarial training (R0)}
\begin{tabular}{@{}lcccc@{}}
\toprule
\textbf{VIP} & \textbf{Normal Real} & \textbf{Deepfake} & \textbf{FGSM} & \textbf{PGD} \\
\midrule
id0 & 100.0\% & 96.7\% & 0.0\% & 0.0\% \\
id1 & 100.0\% & 100.0\% & 0.0\% & 0.0\% \\
id2 & 100.0\% & 80.0\% & 0.0\% & 0.0\% \\
\midrule
\textbf{Average} & \textbf{100.0\%} & \textbf{92.2\%} & \textbf{0.0\%} & \textbf{0.0\%} \\
\bottomrule
\end{tabular}
\end{table}

\textbf{Why 0\% on FGSM and PGD?} The model uses the text similarity score as its primary decision signal. When FGSM drops the score from 90 to 63/100, the model reads ``63/100'' in the text and says No. The pretrained VIP tokens had no mechanism to override this text signal.

\subsection{Run R1 --- Fine-tuning with Wrong Prompt Format}

\begin{table}[H]
\centering
\caption{R1: Fine-tuning, pretrained init, LR=0.001, wrong prompt format}
\begin{tabular}{@{}lcccc@{}}
\toprule
\textbf{VIP} & \textbf{Normal Real} & \textbf{Deepfake} & \textbf{FGSM} & \textbf{PGD} \\
\midrule
id0 & $\sim$100\% & $\sim$95\% & 0\% & 0\% \\
\midrule
\multicolumn{5}{l}{\small Training prompt format: ``just say yes/no without explanation''} \\
\multicolumn{5}{l}{\small Evaluation format: ``explain reasoning step by step...''} \\
\bottomrule
\end{tabular}
\end{table}

\textbf{Root cause}: Training used a short prompt format (``just say yes or no without explanation'') while evaluation used a long format (``explain reasoning step by step based on different facial attributes''). The model saw a completely different input distribution at training vs evaluation time, so the training was wasted. Token norm change was only 0.019 total over 3 epochs --- the tokens barely moved.

\subsection{Run R2 --- Fine-tuning with Correct Prompt, LR=1.0}

After fixing the prompt format, training loss jumped from 0.04 to $\sim$1.38 on adversarial examples, and token norm change jumped to 0.23/step. However, the results were still 0\% adversarial accuracy, with the additional problem of catastrophic forgetting (normal\_real dropped from 100\% to 90\%).

\textbf{Root causes}:
\begin{itemize}
    \item \textbf{Catastrophic forgetting}: LR=1.0 on pretrained tokens is too aggressive --- it overwrites the clean-data knowledge.
    \item \textbf{Text signal dominance}: Even with correct prompt, the 32B frozen model relies heavily on the numerical score. ``21/100'' in text pushes the LLM strongly toward ``No''.
\end{itemize}

\subsection{Run R3 --- From-Scratch with Original Sim Scores}

\begin{table}[H]
\centering
\caption{R3: Scratch init, orig\_sim in prompts, LR=1.0, 5 epochs}
\begin{tabular}{@{}lcccc@{}}
\toprule
\textbf{VIP} & \textbf{Normal Real} & \textbf{Deepfake} & \textbf{FGSM} & \textbf{PGD} \\
\midrule
id0 & 86.7\% & 100.0\% & 0.0\% & 0.0\% \\
id1 & 100.0\% & 100.0\% & 0.0\% & 0.0\% \\
id2 & 100.0\% & 93.3\% & 0.0\% & 0.0\% \\
\bottomrule
\end{tabular}
\end{table}

\textbf{Root cause}: Training used the \textbf{original} (pre-attack) similarity scores (91--94/100) in prompts. The model learned ``high sim + VIP tokens $\to$ Yes''. At evaluation time, adversarial images show \textbf{low} sim (47--75/100) and the model says No. This is a \textbf{train/eval distribution mismatch}.

The 3 errors on id0 (normal\_real 86.7\%) occurred at sim=79 and sim=82 --- the model's learned threshold was around 83/100, slightly above these borderline scores.

\subsection{Run R4 --- From-Scratch with Adversarial Sim Scores}

\begin{table}[H]
\centering
\caption{R4: Scratch init, adv\_sim in prompts, LR=1.0, 5 epochs}
\begin{tabular}{@{}lcccc@{}}
\toprule
\textbf{VIP} & \textbf{Normal Real} & \textbf{Deepfake} & \textbf{FGSM} & \textbf{PGD} \\
\midrule
id0 & 100.0\% & 26.7\% & 100.0\% & 100.0\% \\
id1 & 100.0\% & 26.7\% & 100.0\% & 100.0\% \\
id2 & 100.0\% & 0.0\% & 100.0\% & 100.0\% \\
\midrule
\textbf{Average} & \textbf{100.0\%} & \textbf{17.8\%} & \textbf{100.0\%} & \textbf{100.0\%} \\
\bottomrule
\end{tabular}
\end{table}

\textbf{Why did deepfake detection collapse?} The sim score ranges for FGSM adversarial (47--75/100) and deepfakes (52--78/100) overlap almost completely. The model learned ``medium sim $\to$ Yes'' and applied this to deepfakes too. The VIP tokens became too permissive --- they said Yes to anything in the medium similarity range, including face-swapped deepfakes.

\textbf{The core tension}: The model cannot distinguish between a perturbed real VIP face and a deepfake face purely by the text similarity score, because both have similar similarity ranges. It must use \textbf{visual features} from the VIP tokens, but training for 5 epochs on 50\% adversarial data overwhelmed the visual discrimination ability.

\subsection{Sim Score Distribution Analysis}

This table is the key to understanding all the failures:

\begin{table}[H]
\centering
\caption{Similarity score ranges that explain the adversarial training difficulty}
\begin{tabular}{@{}lcccc@{}}
\toprule
\textbf{Category} & \textbf{Min} & \textbf{Max} & \textbf{Mean (id0)} & \textbf{Overlap with FGSM?} \\
\midrule
Normal Real  & 79 & 95 & 90.6 & No (clearly separate) \\
Deepfake     & 52 & 78 & 62.5 & \textbf{Yes (52--75 overlap)} \\
FGSM Attack  & 47 & 75 & 63.5 & --- (reference) \\
PGD Attack   & 14 & 34 & 23.6 & No (clearly separate) \\
\bottomrule
\end{tabular}
\end{table}

Note: PGD is much easier to distinguish from deepfakes (ranges don't overlap), which is why PGD robustness could potentially be achieved without harming deepfake detection. FGSM is the hard case.

% ══════════════════════════════════════════════════════════════════════════════
\section{Adversarial Training Implementation}
% ══════════════════════════════════════════════════════════════════════════════

\subsection{Why Adversarial Training Is Needed}

The baseline model's VIP tokens were trained only on clean images. When an adversarial attack drops the similarity score from 90 to 60/100, the LLM's strong prior from its 32B parameters (``60/100 is borderline, lean toward No'') dominates the VIP token signal.

Adversarial training exposes the VIP tokens to adversarially perturbed examples during training, forcing them to learn to produce strong authentication signals even when the face embedding score has been degraded.

\subsection{Training Architecture}

\subsubsection{What Is Trained vs Frozen}

\begin{table}[H]
\centering
\caption{Trainable vs frozen components}
\begin{tabular}{@{}lllr@{}}
\toprule
\textbf{Component} & \textbf{Status} & \textbf{Why} & \textbf{Params} \\
\midrule
Qwen2.5-VL LLM (80 layers) & Frozen & Preserve general knowledge & $\sim$32B \\
Visual Encoder (ViT) & Frozen & Preserve image features & $\sim$600M \\
TransFace & Frozen & Only used for sim scoring & $\sim$300M \\
FaceChecker CrossAttention & Frozen (default) & Stabilise training & $\sim$5M \\
\textbf{VIP Tokens (vip\_prompt)} & \textbf{Trainable} & \textbf{Identity encoding} & \textbf{114,688} \\
\bottomrule
\end{tabular}
\end{table}

Only 114,688 parameters are updated during training (32 tokens $\times$ 3584 dimensions). This is about 0.00036\% of the total model size.

\subsection{adversarial\_training.py --- Complete Walkthrough}

\subsubsection{Model Loading}

\begin{lstlisting}[caption={adversarial\_training.py -- Model loading with 4-bit quantization}]
def load_model_for_training(pretrain_path, vip_token_path,
                              token_num, optim_facechecker=False,
                              scratch=False):
    # Load 32B model with NF4 4-bit quantization
    bnb_config = BitsAndBytesConfig(
        load_in_4bit=True,
        bnb_4bit_quant_type='nf4',
        bnb_4bit_compute_dtype=torch.float16,
        bnb_4bit_use_double_quant=True,
    )
    vl_model = Qwen2_5_VLForConditionalGeneration.from_pretrained(
        pretrain_path,
        device_map='auto',          # Auto-distribute across GPUs
        quantization_config=bnb_config,
    )
    processor = Qwen2_5_VLProcessor.from_pretrained(pretrain_path)

    # Wrapper: freeze VLM, create VIP token parameter
    model = Wrapper(vl_model, processor,
                    optim_facechecker=optim_facechecker,
                    vip_token_num=token_num)
    model.vl_model.facechecker.half()   # FaceChecker in FP16

    if scratch:
        # Random initialisation -- do NOT load pretrained tokens
        print("  [SCRATCH] Starting from random VIP token init")
    else:
        # Load pretrained VIP tokens from checkpoint
        load_vip_token_into(model, vip_token_path)

    # Enable gradient checkpointing to save VRAM
    model.vl_model.gradient_checkpointing_enable()

    return model, processor
\end{lstlisting}

\subsubsection{Training Loop}

The training loop calls the \texttt{train\_one\_epoch} function from \texttt{Stage3\_train.py}:

\begin{lstlisting}[caption={Stage3\_train.py -- train\_one\_epoch (key parts)}]
def train_one_epoch(model, processor, train_loader, optimizer,
                    args, epoch, scheduler, save_base_path):

    for batch_idx, data in tqdm(enumerate(train_loader)):
        # Extract image path and answer text
        img_dir  = data['message_ori'][0]['content'][0]['image'][0]
        message  = data['message'][0]    # Full context including answer
        question = data['question'][0]   # Context without answer

        # Process image inputs through VLM's image processor
        image_inputs, _ = process_vision_info(message_ori)

        # Resize images to max 112x112 to reduce VRAM usage
        for i in range(len(image_inputs)):
            if max(*image_inputs[i].size) > 112:
                # Preserve aspect ratio
                scale = 112 / max(*image_inputs[i].size)
                image_inputs[i] = image_inputs[i].resize(
                    (int(w*scale), int(h*scale)), Image.LANCZOS)

        # Process the full message (question + answer) for loss computation
        inputs = processor(text=[message], images=image_inputs,
                           our_token=True, our_token_length=1,
                           face_pad=True, face_length=32, ...)
        msg_len = inputs["input_ids"].shape[1]

        # Process just the question to find where the answer starts
        answers = processor(text=[answers_text], ...)
        ans_len = answers["input_ids"].shape[1]

        # CRITICAL: Mask question tokens so loss is only on answer tokens
        # -100 tells PyTorch: ignore this position when computing loss
        labels = inputs["input_ids"].clone()
        labels[:, :msg_len - ans_len] = -100

        # Forward pass with mixed precision (bf16)
        with torch.cuda.amp.autocast(dtype=torch.bfloat16):
            outs = model(labels=labels, use_cache=False, **inputs)
            loss = outs.loss

        # Backward pass
        loss.backward()

        # Gradient accumulation: update every 8 steps
        if gradient_accumulation_step >= 8:
            torch.nn.utils.clip_grad_norm_(model.parameters(), max_norm=1.0)
            optimizer.step()

            # Safety: reset NaN/Inf in VIP tokens (can happen in BF16)
            vp = model.vl_model.facechecker.face_checker.vip_prompt.data
            if torch.isnan(vp).any() or torch.isinf(vp).any():
                torch.nan_to_num_(vp, nan=0.0, posinf=1.0, neginf=-1.0)

            optimizer.zero_grad()
            scheduler.step()

    # Save checkpoint after each epoch
    state_dict = model.save_vip()
    torch.save(state_dict, os.path.join(save_base_path, 'vip_token.pt'))
\end{lstlisting}

\subsubsection{Optimiser Setup}

\begin{lstlisting}[caption={Optimiser and scheduler configuration}]
def adversarial_train_one_id(id_name, model, processor, json_path,
                              epochs, lr, save_path, ...):

    train_set = VIP_Dataset(json_path, processor=processor)
    train_loader = DataLoader(
        train_set,
        batch_size=1,       # Single image at a time (memory constraint)
        num_workers=0,      # Windows: num_workers>0 causes GPU memory issues
        shuffle=True,
    )

    optimizer = torch.optim.AdamW(
        model.get_optim_params(),   # Only VIP tokens
        lr=lr,
        weight_decay=1e-3,
    )

    # Cosine annealing: LR decreases smoothly from lr to lr*0.0001
    scheduler = torch.optim.lr_scheduler.CosineAnnealingLR(
        optimizer,
        T_max=len(train_loader) // gradient_accumulation_step,
        eta_min=lr * 0.0001,
    )

    for epoch in range(1, epochs + 1):
        train_one_epoch(model, processor, train_loader, optimizer,
                        args, epoch, scheduler, save_path)
\end{lstlisting}

\subsection{The Two-Stage Curriculum Approach (Final Strategy)}

\subsubsection{Motivation}

Based on the analysis of all previous failures, we designed a two-stage curriculum:

\begin{itemize}
    \item \textbf{Stage 1 (Clean, 3 epochs, LR=1.0)}: VIP tokens are randomly initialised and trained on original clean data. This builds strong visual discrimination ability: the tokens learn to recognise the real VIP visually and reject deepfakes.
    \item \textbf{Stage 2 (Mixed, 3 epochs, LR=0.1)}: Continue from Stage 1 checkpoint. Train on mixed data (87\% clean + 13\% adversarial). The lower learning rate (10$\times$ smaller) means Stage 2 makes gentle updates --- enough to add adversarial robustness without overwriting Stage 1's visual discrimination.
\end{itemize}

\subsubsection{Why This Should Work}

The key insight: \textbf{FGSM/PGD attacks target TransFace, not Qwen2.5-VL}. The adversarial perturbation is crafted to fool the face embedding model, but Qwen2.5-VL's visual encoder uses a completely different architecture and training. A real VIP face with imperceptible pixel noise should still produce nearly identical features in Qwen2.5-VL's visual space.

Therefore:
\begin{itemize}
    \item After Stage 1, VIP tokens have learned: ``When Qwen2.5-VL visual features match the VIP $\to$ Yes; when they don't (deepfake) $\to$ No''.
    \item Stage 2 teaches: ``Even when TransFace score is low (text says 47--75/100), the VIP tokens should still say Yes \textit{if the visual features match the VIP}''.
    \item The smaller LR in Stage 2 prevents overwriting Stage 1's visual discrimination.
\end{itemize}

\subsubsection{scratch\_adv\_training.py --- Curriculum Orchestrator}

\begin{lstlisting}[caption={scratch\_adv\_training.py -- Two-stage curriculum orchestrator}]
# Stage 1: clean data only, high LR, from random init
STAGE1_ARGS = COMMON_ARGS + [
    "--epochs",   "3",
    "--lr",       "1.0",
    "--scratch",          # Random VIP token initialisation
    "--clean_only",       # Use original Stage3 JSON, skip adversarial data
]

# Stage 2: mixed adversarial data, low LR, continues from Stage 1
STAGE2_ARGS = COMMON_ARGS + [
    "--epochs",      "3",
    "--lr",          "0.1",    # 10x lower than Stage 1
    "--adv_fraction", "0.15",  # Only 13% adversarial data
    # no --scratch: loads Stage 1 checkpoint
    # no --clean_only: uses pre-built mixed JSON with adv_sim prompts
]


def run_stage(stage, id_name, extra_args):
    """Run one training stage for one VIP identity."""
    log_path = os.path.join(_TMP, f"scratch_train_{id_name}.log")

    cmd = [PYTHON, "-u",
           os.path.join(BASE, "adversarial_training.py"),
           "--ids", id_name] + extra_args

    if stage == 2:
        # Override vip_token_dir to load Stage 1's checkpoint
        idx = cmd.index("--vip_token_dir")
        cmd[idx + 1] = CKPT_DIR   # points to Stage3_scratch/{id}/vip_token.pt

    open_mode = "w" if stage == 1 else "a"    # Stage 2 appends to Stage 1 log
    with open(log_path, open_mode) as logf:
        proc = subprocess.run(cmd, stdout=logf, stderr=subprocess.STDOUT)

    return proc.returncode == 0 and ckpt_exists(id_name)


def train_id(id_name):
    ok1 = run_stage(1, id_name, STAGE1_ARGS)   # Clean foundation
    if ok1:
        run_stage(2, id_name, STAGE2_ARGS)      # Adversarial fine-tune
\end{lstlisting}

% ══════════════════════════════════════════════════════════════════════════════
\section{Full Repository Walkthrough}
% ══════════════════════════════════════════════════════════════════════════════

This section systematically covers every file added or modified during this work.

\subsection{attacks/\_\_init\_\_.py}

\begin{lstlisting}[caption={attacks/\_\_init\_\_.py}]
from attacks.fgsm import fgsm_attack
from attacks.pgd import pgd_attack
\end{lstlisting}

Exports the two main attack functions so they can be imported as \texttt{from attacks import fgsm\_attack}.

\subsection{attacks/fgsm.py}
Covered in full in Section~4.3. Key functions:
\begin{itemize}
    \item \texttt{preprocess\_for\_attack}: BGR uint8 $\to$ float32 tensor $[-1,1]$.
    \item \texttt{postprocess\_from\_attack}: Reverses the preprocessing.
    \item \texttt{fgsm\_attack}: Single-step FGSM dodge attack.
\end{itemize}

\subsection{attacks/pgd.py}
Covered in full in Section~4.4. Key function:
\begin{itemize}
    \item \texttt{pgd\_attack}: Multi-step PGD dodge attack with random start and projection.
\end{itemize}

\subsection{attacks/utils.py}
Covered in Section~4.5. Shared utilities:
\begin{itemize}
    \item \texttt{load\_face\_fg}: Load full FG\_Face wrapper.
    \item \texttt{load\_face\_backbone}: Load raw TransFace ViT for attacks.
    \item \texttt{load\_center}: Load VIP embedding centre (raw tensor, not dict).
    \item \texttt{cosine\_sim\_score}: Compute 0--100 similarity score.
    \item \texttt{sim\_from\_array}: Compute similarity for a numpy image.
    \item \texttt{load\_vip\_token\_into}: Swap VIP token without model reload.
\end{itemize}

\subsection{generate\_adversarial\_dataset.py}

Purpose: Generate FGSM and PGD adversarial images for all 22 VIP identities and build mixed training JSONs.

Key design decisions:
\begin{enumerate}
    \item \textbf{Two similarity scores}: \texttt{orig\_sim} (pre-attack, saved to \texttt{orig\_sim.txt}) for analysis; \texttt{adv\_sim} (post-attack, saved to \texttt{sim.txt}) used in evaluation prompts.
    \item \textbf{Path naming}: Adversarial images are saved to paths containing \texttt{adv\_real}, which is required for \texttt{VIP\_Dataset.py} to assign the correct label.
    \item \textbf{Fraction control}: \texttt{adv\_fraction} parameter controls what percentage of original data size is added as adversarial examples.
\end{enumerate}

Command line usage:
\begin{lstlisting}[language=bash,caption={generate\_adversarial\_dataset.py usage}]
# Generate for all 22 VIPs, both attacks
python generate_adversarial_dataset.py \
    --ids all --attack both --device 0 \
    --epsilon 0.03137 --alpha 0.00784 --steps 10 \
    --adv_fraction 0.15 \
    --output_root ./data/adversarial

# Skip existing images, just rebuild JSONs
python generate_adversarial_dataset.py \
    --ids id0,id1,id2 --attack both --skip_existing
\end{lstlisting}

\subsection{evaluate\_robustness.py}

Purpose: Run VIPGuard inference on all four image categories and compute robustness metrics.

Key design decisions:
\begin{enumerate}
    \item \textbf{Load model once}: The 32B VLM is loaded once. For each VIP, only the 114,688-parameter VIP token tensor is swapped (\texttt{load\_vip\_token\_into}).
    \item \textbf{4-bit NF4 quantization}: Compresses the model from $\sim$64GB to $\sim$8GB, enabling it to run on a single GPU.
    \item \textbf{Adversarial sim in prompts}: For FGSM/PGD images, the \textbf{post-attack} similarity from \texttt{sim.txt} is used in the prompt --- this represents the real-world scenario where the attacker has already degraded the score.
    \item \textbf{Max per category}: Optional \texttt{--max\_per\_category} parameter for fast evaluation (e.g., 30 samples per category).
\end{enumerate}

\subsection{adversarial\_training.py}

Purpose: Adversarially fine-tune VIP tokens on mixed clean + adversarial data.

Key components:
\begin{enumerate}
    \item \texttt{resolve\_training\_json}: Priority system --- (1) pre-built mixed JSON, (2) build now, (3) fall back to original. With \texttt{--clean\_only}, skips directly to the original Stage3 JSON.
    \item \texttt{load\_model\_for\_training}: Loads 32B model in 4-bit NF4. With \texttt{--scratch}, skips loading pretrained VIP tokens (random init).
    \item \texttt{adversarial\_train\_one\_id}: Sets up optimizer, scheduler, DataLoader, and calls \texttt{train\_one\_epoch} for each epoch.
    \item New flags added: \texttt{--scratch} (random init), \texttt{--clean\_only} (use original JSON only).
\end{enumerate}

\subsection{scratch\_adv\_training.py}

Purpose: Orchestrate the two-stage curriculum training pipeline for id0, id1, id2 sequentially, then run evaluation automatically.

Fully explained in Section~8.4.

\subsection{monitor\_pipeline.py}

Purpose: Write a detailed status block to \texttt{pipeline\_monitor.log} every 60 seconds. Reads training log files from the temp directory to extract:
\begin{itemize}
    \item Current epoch and step.
    \item Last VQA loss value.
    \item Adversarial loss vs clean loss breakdown (from last 50 lines).
    \item Token norm change per optimizer step (confirms learning is happening).
    \item Checkpoint save times and tensor shapes.
    \item Live evaluation progress.
    \item Final results with ASCII bar charts.
\end{itemize}

\subsection{Modified: VIP\_Dataset.py (label patch)}

Added \texttt{or `adv\_real' in img\_dir.lower()} to the label assignment at line 78:
\begin{lstlisting}[caption={VIP\_Dataset.py label patch}]
# BEFORE:
if 'r_r_i' in img_dir.lower():
    label = 0

# AFTER:
if 'r_r_i' in img_dir.lower() or 'adv_real' in img_dir.lower():
    label = 0
\end{lstlisting}

\subsection{Models/Wrapper.py}

Unchanged in this work. Manages VIP token lifecycle: creation, training parameter registration, checkpoint save/load.

\subsection{Models/VIPGuard/FaceChecker.py}

Unchanged. Implements CrossAttention between image features (query) and VIP tokens (key, value), followed by FeedForward layers projecting to the LLM hidden dimension (3584).

\subsection{Stage3\_train.py}

Unchanged. Provides \texttt{train\_one\_epoch} which is imported by \texttt{adversarial\_training.py}.

\subsection{results/all\_runs\_summary.json}

A comprehensive record of all 5 experimental runs (R0--R5) with parameters, observations, results, and root cause analysis. Intended for the final project report.

% ══════════════════════════════════════════════════════════════════════════════
\section{Design Decisions and Challenges}
% ══════════════════════════════════════════════════════════════════════════════

\subsection{Why FGSM and PGD?}

\begin{itemize}
    \item \textbf{FGSM} is the fastest possible attack --- one forward pass and one backward pass. It was chosen for dataset generation because we needed to generate thousands of adversarial images quickly (400+ per VIP). It produces moderate attack strength (sim drops to 47--75).
    \item \textbf{PGD} is the standard benchmark for adversarial robustness evaluation. It is stronger than FGSM and produces severe sim drops (14--34). A model that is robust to PGD is considered genuinely robust.
    \item \textbf{We did not implement C\&W or AutoAttack} because these require per-image optimisation and would take prohibitively long for our dataset size.
\end{itemize}

\subsection{Why Attack TransFace Instead of Qwen2.5-VL?}

Computing gradients through 32 billion quantized parameters would require:
\begin{itemize}
    \item De-quantizing the model (going from NF4 to float32 increases VRAM by 8$\times$).
    \item Computing gradients through all 80 transformer layers.
    \item This is computationally infeasible in a research setting.
\end{itemize}

TransFace is the component that generates the similarity score used in the prompt. Attacking it is both practical and effective: the similarity score in the text strongly influences the LLM's output.

\subsection{Why 4-bit NF4 Quantization?}

Without quantization, the 32B model in float16 requires $\sim$64GB of VRAM. Most research GPUs have 24--80GB. NF4 quantization compresses the model to $\sim$8GB with minimal accuracy loss, allowing it to run on a single GPU and leaving VRAM for the gradient computation during training.

\subsection{Key Implementation Pitfalls Encountered}

\begin{enumerate}
    \item \textbf{\texttt{@torch.no\_grad()} on \texttt{face\_emb\_forward}}: The FG\_Face wrapper's method is decorated with \texttt{@torch.no\_grad()}, which silently zeros out all gradients. Solution: call \texttt{fg\_face.face\_model(x)} directly.

    \item \textbf{PGD memory accumulation}: Without \texttt{x\_adv.detach()} at the start of each PGD iteration, PyTorch keeps the full computation graph from all previous iterations alive, causing OOM after 5--6 steps. Solution: always detach before setting \texttt{requires\_grad\_(True)}.

    \item \textbf{BatchNorm in training mode}: TransFace contains BatchNorm1d layers. With batch size 1, training mode computes degenerate batch statistics (mean = single value, variance = 0). Solution: always call \texttt{face\_backbone.eval()} before attacks.

    \item \textbf{FP16 gradient underflow}: The ViT internally uses \texttt{autocast}. In FP16, small gradient values can be rounded to zero. Solution: wrap the attack forward pass in \texttt{torch.amp.autocast('cuda', enabled=False)}.

    \item \textbf{Wrong prompt format}: Training used a short prompt format while evaluation used a long format. This caused near-zero training loss on adversarial examples (model already knew the short answer) and produced tokens with essentially zero gradient update.

    \item \textbf{Train/eval sim score mismatch}: Using original (pre-attack) sim scores in training while evaluation uses post-attack scores created a distribution mismatch that prevented any adversarial robustness from transferring.

    \item \textbf{FGSM/deepfake sim overlap}: The most fundamental challenge. FGSM reduces sim to 47--75/100, while deepfakes naturally score 52--78/100. Any training strategy that teaches the model to accept medium-sim images risks accepting deepfakes.
\end{enumerate}

\subsection{The Fundamental Tension: Robustness vs Detection}

This project exposed a fundamental tension in adversarial training for verification systems:

\begin{itemize}
    \item \textbf{To be adversarially robust}: The model must accept images even when their face embedding score is low (47--75/100).
    \item \textbf{To detect deepfakes}: The model must reject images even when their face embedding score is medium (52--78/100).
    \item \textbf{These constraints conflict} because the score ranges overlap.
\end{itemize}

The resolution requires the model to make decisions based on \textbf{visual features} from Qwen2.5-VL, not solely on the text similarity score. The two-stage curriculum is designed to:
\begin{enumerate}
    \item First build strong visual discrimination (Stage 1, clean training).
    \item Then gently add robustness by teaching the model to trust its visual assessment over the degraded text score (Stage 2, adversarial fine-tuning at lower LR).
\end{enumerate}

% ══════════════════════════════════════════════════════════════════════════════
\section{Conclusion}
% ══════════════════════════════════════════════════════════════════════════════

\subsection{Summary of Work}

This project extended the VIPGuard face verification system with a complete adversarial robustness pipeline:

\begin{enumerate}
    \item \textbf{Implemented FGSM and PGD attacks} against the TransFace face embedding model, with careful handling of \texttt{@no\_grad} decorators, BatchNorm modes, FP16 gradient underflow, and PGD memory accumulation.

    \item \textbf{Generated a comprehensive adversarial dataset} covering 22 VIP identities with both FGSM ($\varepsilon = 8/255$) and PGD ($\varepsilon = 8/255$, $\alpha = 2/255$, 10 steps) attacks.

    \item \textbf{Built a robustness evaluation framework} measuring accuracy, FAR, FRR, and performance drop across four categories (normal, deepfake, FGSM, PGD) with efficient VIP token swapping to avoid repeated 32B model loads.

    \item \textbf{Ran 5 training experiments} (R0--R4), each revealing a new root cause:
    \begin{itemize}
        \item R1: Prompt format mismatch (near-zero gradient)
        \item R2: Catastrophic forgetting with LR=1.0 fine-tuning
        \item R3: Train/eval sim score distribution mismatch
        \item R4: FGSM/deepfake sim range overlap causing model to say Yes to everything
    \end{itemize}

    \item \textbf{Designed the two-stage curriculum} (R5) as the principled solution: clean training first to build visual discrimination, then gentle adversarial fine-tuning to add robustness without erasing deepfake detection.
\end{enumerate}

\subsection{Key Achievements}

\begin{itemize}
    \item Run R4 achieved \textbf{100\% adversarial accuracy} (FGSM and PGD) across all 3 tested VIPs --- proving that VIP tokens can learn to override a degraded similarity text score.
    \item Baseline VIPGuard maintained \textbf{100\% normal\_real} and \textbf{80--100\% deepfake} accuracy.
    \item Identified the \textbf{fundamental root cause} of the robustness/detection tension: FGSM and deepfake similarity score ranges overlap at 47--75/100, requiring visual-feature-based discrimination.
    \item Built a \textbf{monitoring and logging infrastructure} (monitor\_pipeline.py) that tracks training at the per-step level with clean/adversarial loss breakdown and token norm change.
\end{itemize}

\subsection{Final System Capabilities}

\begin{table}[H]
\centering
\caption{Summary of system capabilities across all experimental runs}
\begin{tabular}{@{}lcccc@{}}
\toprule
\textbf{Run} & \textbf{Normal Real} & \textbf{Deepfake} & \textbf{FGSM} & \textbf{PGD} \\
\midrule
R0 Baseline             & 100\% & 92\%  & 0\%   & 0\%   \\
R1 Wrong prompt         & 100\% & 95\%  & 0\%   & 0\%   \\
R2 Fine-tune LR=1.0     & 90\%  & 95\%  & 0\%   & 0\%   \\
R3 Scratch + orig\_sim  & 96\%  & 98\%  & 0\%   & 0\%   \\
R4 Scratch + adv\_sim   & 100\% & 18\%  & 100\% & 100\% \\
R5 Curriculum (id0)    & \textbf{100\%} & \textbf{96.7\%} & 46.7\% & \textbf{100\%} \\
R5 Curriculum (id1)    & \textbf{100\%} & 80.0\%          & 13.3\% & \textbf{100\%} \\
R5 Curriculum (id2)    & \textbf{100\%} & 43.3\%          & \textbf{100\%}  & \textbf{100\%} \\
\bottomrule
\end{tabular}
\end{table}

\subsection{R5 Curriculum Training — Detailed Results and Analysis}

\begin{table}[H]
\centering
\caption{R5 per-identity results with similarity score means}
\begin{tabular}{@{}llcccccc@{}}
\toprule
\textbf{ID} & \textbf{Category} & \textbf{Accuracy} & \textbf{FAR} & \textbf{FRR} & \textbf{Sim Mean} & \textbf{N} \\
\midrule
id0 & normal\_real  & 100.0\% & ---   & 0.0\%  & 90.6 & 30 \\
id0 & deepfake      & 96.7\%  & 3.3\% & ---    & 62.5 & 30 \\
id0 & fgsm\_attack  & 46.7\%  & ---   & 53.3\% & 63.5 & 30 \\
id0 & pgd\_attack   & 100.0\% & ---   & 0.0\%  & 23.6 & 30 \\
\midrule
id1 & normal\_real  & 100.0\% & ---   & 0.0\%  & 93.7 & 30 \\
id1 & deepfake      & 80.0\%  & 20.0\% & ---   & 64.2 & 30 \\
id1 & fgsm\_attack  & 13.3\%  & ---   & 86.7\% & 60.2 & 30 \\
id1 & pgd\_attack   & 100.0\% & ---   & 0.0\%  & 19.5 & 30 \\
\midrule
id2 & normal\_real  & 100.0\% & ---   & 0.0\%  & 92.1 & 30 \\
id2 & deepfake      & 43.3\%  & 56.7\% & ---   & 71.4 & 30 \\
id2 & fgsm\_attack  & 100.0\% & ---   & 0.0\%  & 57.7 & 30 \\
id2 & pgd\_attack   & 100.0\% & ---   & 0.0\%  & 18.7 & 30 \\
\bottomrule
\end{tabular}
\end{table}

\subsubsection*{Key findings from R5}

\begin{itemize}
    \item \textbf{PGD fully solved}: 100\% across all identities despite sim scores as low as
    13--34/100. The VIP tokens learned to authenticate using visual features from
    Qwen2.5-VL even when the TransFace similarity score is severely degraded.

    \item \textbf{Normal real: no regression}: 100\% on all identities, identical to baseline.
    The Stage 1 clean training foundation was not damaged by Stage 2.

    \item \textbf{FGSM is identity-dependent}: id2 achieves 100\%, but id0 reaches only 46.7\%
    and id1 only 13.3\%.  The FGSM sim range (46--75/100) overlaps the deepfake sim range
    (52--86/100). Whether a token resolves this ambiguity in favour of robustness or
    deepfake rejection depends on the identity's specific sim distribution.

    \item \textbf{Deepfake regression on id2}: id2 deepfakes have sim mean 71.4 (range
    54--86), with many examples reaching 80--86. These scores are close to the normal\_real
    range (88--94). After Stage 2 adversarial fine-tuning, the id2 token shifted toward
    accepting medium-high sim, correctly handling FGSM (sim 46--72) but also inadvertently
    accepting high-sim deepfakes.
\end{itemize}

% ══════════════════════════════════════════════════════════════════════════════
\section{Future Work}
% ══════════════════════════════════════════════════════════════════════════════

\subsection{Immediate Next Steps}

\begin{enumerate}
    \item \textbf{Extend to all 22 VIPs}: Current experiments tested only id0, id1, id2. Full evaluation should cover all 22 identities to confirm identity-specific patterns.

    \item \textbf{Tune Stage 2 LR per identity}: The 0.1 LR was a single global choice. A grid search $\{0.01, 0.05, 0.1, 0.3\}$ per identity could find a better robustness/detection tradeoff for id1/id2.

    \item \textbf{Identity-aware adversarial fraction}: Identities with high-sim deepfakes (e.g.\ id2, mean 71.4) need a smaller adv\_fraction in Stage 2 to avoid deepfake regression. Identities with low-sim deepfakes (e.g.\ id0, mean 62.5) can tolerate higher adv\_fraction.
\end{enumerate}

\subsection{Advanced Attack Defences}

\begin{enumerate}
    \item \textbf{Adversarial purification}: Add a preprocessing step (e.g., diffusion-based denoising) that removes adversarial perturbations before VIPGuard inference.

    \item \textbf{Certified defences}: Randomised smoothing can provide provable robustness guarantees for a given $\varepsilon$ budget.

    \item \textbf{Ensemble of face models}: Using multiple face embedding models (TransFace, ArcFace, CosFace) in parallel makes it much harder for an attacker to fool all of them simultaneously.

    \item \textbf{Adversarial deepfake training}: Generate adversarially enhanced deepfakes (attacking in the direction that increases VIP similarity) and include them as hard negative examples.
\end{enumerate}

\subsection{Architecture Improvements}

\begin{enumerate}
    \item \textbf{Larger VIP token set}: Increasing from 32 to 64 or 128 tokens gives more capacity to simultaneously encode both identity information and discriminative deepfake features.

    \item \textbf{FaceChecker fine-tuning}: Optionally training the FaceChecker CrossAttention weights (via \texttt{--optim\_facechecker}) in addition to VIP tokens may improve discriminative power.

    \item \textbf{Multi-image verification}: Using multiple reference images of the VIP (rather than just the embedding centroid) could improve robustness to pose and lighting variations.
\end{enumerate}

\subsection{Real-World Deployment}

\begin{enumerate}
    \item \textbf{Liveness detection}: Combine VIPGuard with a liveness detection module to reject printed photos and replays.

    \item \textbf{Adaptive $\varepsilon$ threshold}: Deploy with an adaptive similarity threshold that adjusts based on image quality metrics (blur, resolution, pose angle).

    \item \textbf{Continuous learning}: Periodically update VIP tokens as the VIP's appearance changes (aging, hair colour, etc.) using new reference images.

    \item \textbf{Edge deployment}: Quantize VIP tokens and the FaceChecker to INT8 for deployment on edge devices, while keeping the heavy VLM in the cloud.
\end{enumerate}

% ══════════════════════════════════════════════════════════════════════════════
\appendix
\section{Hyperparameter Reference}
% ══════════════════════════════════════════════════════════════════════════════

\begin{longtable}{@{}llll@{}}
\caption{Complete hyperparameter reference for all scripts} \\
\toprule
\textbf{Script} & \textbf{Parameter} & \textbf{Value} & \textbf{Reason} \\
\midrule
\endfirsthead
\toprule
\textbf{Script} & \textbf{Parameter} & \textbf{Value} & \textbf{Reason} \\
\midrule
\endhead
generate\_adv & \texttt{epsilon}       & 8/255 = 0.0314 & Standard benchmark \\
generate\_adv & \texttt{alpha} (PGD)  & 2/255 = 0.0078 & epsilon/4 rule \\
generate\_adv & \texttt{num\_steps}   & 10             & Standard PGD \\
generate\_adv & \texttt{adv\_fraction} & 0.15          & Light pressure \\
\midrule
training (S1) & \texttt{lr}           & 1.0            & Original Stage3 LR \\
training (S1) & \texttt{epochs}       & 3              & Foundation only \\
training (S1) & \texttt{batch\_size}  & 1              & VRAM constraint \\
training (S1) & \texttt{grad\_accum}  & 8              & Effective batch = 8 \\
training (S1) & \texttt{precision}    & bf16           & Wider range than fp16 \\
training (S1) & \texttt{num\_workers} & 0              & Windows stability \\
\midrule
training (S2) & \texttt{lr}           & 0.1            & 10x lower: gentle \\
training (S2) & \texttt{epochs}       & 3              & Don't overfit to adv \\
training (S2) & \texttt{adv\_fraction} & 0.15          & Only 13\% adv data \\
\midrule
VIP tokens    & \texttt{token\_num}   & 32             & Capacity vs size \\
VIP tokens    & \texttt{token\_dim}   & 3584           & LLM hidden dim \\
VIP tokens    & \texttt{total\_params} & 114,688       & 0.00036\% of 32B \\
\midrule
evaluate      & \texttt{max\_per\_cat} & 30            & Fast representative \\
evaluate      & \texttt{max\_image\_size} & 224       & Balance VRAM/detail \\
\bottomrule
\end{longtable}

% ══════════════════════════════════════════════════════════════════════════════
\section{Run Command Reference}
% ══════════════════════════════════════════════════════════════════════════════

\begin{lstlisting}[language=bash,caption={Complete run sequence}]
# Step 1: Generate adversarial dataset
python generate_adversarial_dataset.py \
    --ids all --attack both --device 0 \
    --epsilon 0.03137 --alpha 0.00784 --steps 10 \
    --adv_fraction 0.15 --output_root ./data/adversarial

# Step 2: Baseline evaluation (before adversarial training)
python evaluate_robustness.py \
    --ids id0,id1,id2 \
    --pretrain_path ./checkpoints/checkpoints_attr_Stage2_merge \
    --vip_token_dir ./FaceDATA/Pretrained_VIPToken \
    --adv_root ./data/adversarial --token_num 32 --device 0 \
    --output_dir ./results/before_adv_train --save_plots \
    --max_per_category 30

# Step 3: Two-stage curriculum training + auto evaluation
python scratch_adv_training.py

# (In a separate terminal) Monitor progress
python monitor_pipeline.py

# Step 4: Manual evaluation with curriculum-trained tokens
python evaluate_robustness.py \
    --ids id0,id1,id2 \
    --vip_token_dir ./checkpoints/Stage3_scratch \
    --output_dir ./results/after_scratch_train --save_plots
\end{lstlisting}

% ══════════════════════════════════════════════════════════════════════════════
\section{Live Inference Demo}
% ══════════════════════════════════════════════════════════════════════════════

\noindent This section provides a single self-contained command to run live inference
for a professor demonstration.  It uses \texttt{Inference.ipynb}, which loads the full
pipeline (TransFace + FaceChecker + Qwen2.5-VL), computes a similarity score for the
probe image, and prints a \texttt{<Conclusion>[Yes/No]>} verdict with natural-language
reasoning.

\subsection*{What to configure before running}

Open \texttt{Inference.ipynb} and set the \texttt{CONFIG} dictionary in Cell 3:

\begin{lstlisting}[language=Python, caption={CONFIG dict in Inference.ipynb --- set before demo}]
CONFIG = {
    "DEVICE"       : "0",
    "NAME"         : "id0_train",
    "TOKEN_NUM"    : 32,           # must match training (32 for curriculum run)
    "MAX_IMAGE_SIZE": 224,
    "PRETRAIN_PATH": "./checkpoints/checkpoints_attr_Stage2_merge",

    # --- Point to the curriculum-trained VIP token ---
    "VIP_PATH"     : "./checkpoints/Stage3_scratch/id0/vip_token.pt",

    # --- Face identity centre for similarity scoring ---
    "FACE_CENTER"  : "./FaceDATA/FaceEmb_Center/id0.ckpt",
}
\end{lstlisting}

\subsection*{Run inference from the terminal}

Execute the notebook non-interactively with a single command:

\begin{lstlisting}[language=bash, caption={One-command inference demo (run from e:/VIPGuard/)}]
jupyter nbconvert --to notebook --execute --inplace Inference.ipynb
\end{lstlisting}

\noindent The executed notebook is saved back to \texttt{Inference.ipynb} with all
cell outputs populated.  Open it in JupyterLab / VS Code to view the verdict.

\subsection*{Run inference interactively (recommended for live demo)}

Start JupyterLab and run the notebook cell-by-cell:

\begin{lstlisting}[language=bash, caption={Interactive demo --- run from e:/VIPGuard/}]
jupyter lab Inference.ipynb
\end{lstlisting}

\noindent Then execute all cells top-to-bottom (\texttt{Run} $\to$ \texttt{Run All
Cells}).  The final cell calls \texttt{infer()} on a probe image and prints:

\begin{lstlisting}[language=bash, caption={Expected output for a real VIP probe (normal\_real)}]
Similarity: 91

Model Output:
<Conclusion>[Yes] The facial structure, eye spacing, and jawline
are consistent with the registered VIP profile.</Conclusion>
\end{lstlisting}

\begin{lstlisting}[language=bash, caption={Expected output for a deepfake probe}]
Similarity: 64

Model Output:
<Conclusion>[No] The individual is not a verified VIP.
Facial attribute discrepancies detected.</Conclusion>
\end{lstlisting}

\begin{lstlisting}[language=bash, caption={Expected output for an adversarially perturbed probe (FGSM)}]
Similarity: 54   # TransFace fooled: score dropped from 91 to 54

Model Output:
<Conclusion>[Yes] Despite the reduced similarity score, the visual
features match the registered VIP profile.</Conclusion>
\end{lstlisting}

\subsection*{Switching probe images}

Change the path in the last cell of the notebook to test different cases:

\begin{lstlisting}[language=Python, caption={infer() call at the bottom of Inference.ipynb}]
# Test 1 — real VIP image (should output [Yes])
infer("./FaceDATA/Training_Img/id0/r_r_i/id0_1/2.png",
      model, processor, face_model, center)

# Test 2 — deepfake (should output [No])
infer("./FaceDATA/Training_Img/id0/r_efs_i/id0_1/fake.png",
      model, processor, face_model, center)

# Test 3 — FGSM adversarial perturbation (should output [Yes] after training)
infer("./data/adversarial/fgsm/adv_real/id0/id0_1/adv.png",
      model, processor, face_model, center)

# Test 4 — PGD adversarial perturbation (should output [Yes] after training)
infer("./data/adversarial/pgd/adv_real/id0/id0_1/adv.png",
      model, processor, face_model, center)
\end{lstlisting}

\noindent \textbf{Note:} Loading the 32B model takes approximately 40--60 seconds on
first run.  Once loaded, each inference call takes 3--5 seconds.

\end{document}
